\section{Related work}\label{sec:related}

\paragraph{Statistical forecasting of the Court.}
The modern line begins with the Supreme Court Forecasting Project:
\citet{ruger2004} built a conditional-choice model with input from
political science theory and, over a single term, predicted 75\% of the
Court's affirm/reverse outcomes---against 59.1\% for a panel of
legal experts. The lesson was bruising for doctrine: a small
regression with ideological covariates beat trained lawyers reading
the briefs. A decade later, \citet{katz2017} scaled the exercise with
a random forest over 628 features and nearly two centuries of
decisions (1816--2013), reporting 70.2\% accuracy at the case level
and 71.9\% at the individual vote level, and concluding that the
Court behaves in a way that is, in their words, remarkably consistent
with a machine-learning account. Both efforts, however, evaluate
\emph{in sample across history}: the models see several centuries ---
including eras with much higher unanimity rates---and predict votes
whose contemporaneous context is far removed from today's Court.

\paragraph{The attitudinal model.}
Political science supplies the theoretical prior for why such models
work. \citet{segal2002} argue that justices decide largely as
attitudes wrapped in legal reasoning: revealed ideology dominates
doctrine, and case facts matter mostly as triggers for those
attitudes. Methodologically, \citet{martin2002} estimate dynamic
ideal points for every justice since 1953, giving the field a
continuous measure of the very covariate that our strongest baseline
--- per-justice modal direction---approximates with a single bit per
justice. Our results in Section~\ref{sec:results} are, in effect, a
modern, pre-registered replication of the attitudinal prior on a
nine-term window: the ideology bit is the signal; the structured case
context is not.

\paragraph{Authorship and style.}
The persona condition is a stylometry question in modern dress.
\citet{mosteller1964} settled the authorship of the disputed
\emph{Federalist} papers with function-word frequencies; the field
they founded has since shown that individual writing style is
persistent, measurable, and hard to disguise~\citep{juola2006}. If a
justice's voice is a stable individual signature, fine-tuning a
language model on their opinions should transfer to their future
opinions; the open question---the one this experiment isolates ---
is whether that stylistic footprint carries \emph{decision-relevant}
information beyond the ideology bit, or merely rhetorical habits.

\paragraph{Language models and efficient fine-tuning.}
The tooling for the persona condition is recent. QLoRA
\citep{dettmers2023} makes full-model fine-tuning tractable on free
single-GPU tiers by quantizing the base model and training low-rank
adapters, and open-weight families such as Llama~3
\citep{llama3} provide a competent base. This is what converts the
persona question from an institutional project into a zero-budget
one: nine adapters, one per justice, trained on free-tier compute,
each sealed alongside the data that produced it.

\paragraph{Positioning.}
Against this background, the present project narrows the claim in
three ways that we believe matter for honesty of the result. The
window is short and modern (OT2015--OT2023, an era of narrow
majorities), so historical unanimity cannot inflate accuracy. The
evaluation is pre-registered and forward-holding: a sealed set of
fifty 5--4 decisions, chosen deterministically before any model
touched them, is the only arbiter of the final conditions. And the
bar is fixed \emph{before} the challengers run, at the vote level,
where the attitudinal bit is strongest---63.7\% on our window ---
so that any reported improvement is an improvement over ideology
itself, not over a straw man.
